\documentclass[12pt]{article}


\usepackage{fullpage}
\usepackage{mdframed}
\usepackage{colonequals}
\usepackage{algpseudocode}
\usepackage{algorithm}
\usepackage{tcolorbox}
\usepackage[all]{xy}
\usepackage{proof}
\usepackage{mathtools}
\usepackage{bbm}
\usepackage{amssymb}
\usepackage{amsthm}
\usepackage{amsmath}
\usepackage{amsxtra}
\usepackage{enumitem}
\usepackage{euler}



% Theorem environments
\newtheorem{theorem}{Theorem}[section]
\newtheorem{theorem*}{Theorem}
\newtheorem{definition}[theorem]{Definition}
\newtheorem{corollary}{Corollary}[theorem]
\newtheorem{lemma}[theorem]{Lemma}
\newtheorem{proposition}[theorem]{Proposition}
\newtheorem*{remark}{Remark}
\newtheorem*{claim}{Claim}
\newtheorem*{example}{Example}
\newtheorem*{warning}{Warning}

\newtheorem*{exercisehelper}{Exercise.}
\newenvironment{exercise}[1]{%
  \IfBlankTF{#1}
    {\renewcommand{\exercisehelper}{\textbf{Exercise} \unskip}}
    {\renewcommand\exercisehelper{\textbf{Exercise #1}}}%
  \exercisehelper
}{\endexercisehelper}

\theoremstyle{remark}
\newtheorem*{solution}{Solution}
\newcommand{\mathcat}[1]{\textup{\textbf{\textsf{#1}}}} % for defined terms

\newenvironment{problem}[1]
{\begin{tcolorbox}\noindent\textbf{Problem #1}.}
{\vskip 6pt \end{tcolorbox}}

\newenvironment{enumalph}
{\begin{enumerate}\renewcommand{\labelenumi}{\textnormal{(\alph{enumi})}}}
{\end{enumerate}}

\newenvironment{enumroman}
{\begin{enumerate}\renewcommand{\labelenumi}{\textnormal{(\roman{enumi})}}}
{\end{enumerate}}

\let\H\relax
\let\P\relax
\let\Pr\relax
\let\d\relax
\let\limsup\relax
\let\1\relax
\let\div\relax


\DeclarePairedDelimiter\ceil{\lceil}{\rceil}
\DeclarePairedDelimiter\floor{\lfloor}{ \rfloor}
\newcommand{\fr}[2]{\left(\frac{#1}{#2}\right)}
\newcommand{\norm}[1]{\left\vert\left\vert#1\right\vert\right\vert}
\newcommand{\dbar}{\overline{\partial}}
\newcommand{\d}{\partial}
\newcommand{\RP}{\mathbb{RP}}
\newcommand{\CP}{\mathbb{CP}}
\newcommand{\tbigwedge}{{\textstyle{\bigwedge}}}


\newcommand{\bbni}{\bigbreak \noindent}
\newcommand{\todo}[1]{\textcolor{red}{\textbf{TODO:} #1}}
\newcommand{\red}[1]{\textcolor{red}{#1}}


% Blackboard Font
\newcommand{\N}{\mathbb N}
\newcommand{\Z}{\mathbb Z}
\newcommand{\Q}{\mathbb Q}
\newcommand{\R}{\mathbb R}
\newcommand{\C}{\mathbb C}
\newcommand{\F}{\mathbb F}
\newcommand{\G}{\mathbb G}
\newcommand{\H}{\mathbb H}
\newcommand{\bb}{\mathbb}

% Calligraphic font
\newcommand{\CC}{\mathcal{C}}
\newcommand{\DD}{\mathcal{D}}
\newcommand{\EE}{\mathcal{E}}
\newcommand{\FF}{\mathcal{F}}
\newcommand{\GG}{\mathcal{G}}
\newcommand{\HH}{\mathcal{H}}
\newcommand{\LL}{\mathcal{L}}
\newcommand{\OO}{\mathcal{O}}
\newcommand{\II}{\mathcal{I}}


% Linear Algebra
\DeclareMathOperator{\GL}{GL}
\DeclareMathOperator{\PGL}{PGL}
\DeclareMathOperator{\Mat}{Mat}
\DeclareMathOperator{\Id}{Id}
\DeclareMathOperator{\Tr}{Tr}
\DeclareMathOperator{\tr}{tr}
\DeclareMathOperator{\Nm}{Nm}
\DeclareMathOperator{\opspan}{span}
\DeclareMathOperator{\rk}{rk}
\DeclareMathOperator{\sgn}{sgn}
\DeclareMathOperator{\Sym}{Sym}
\DeclareMathOperator{\Alt}{Alt}
\DeclareMathOperator{\codim}{codim}



% Category Theory / Homological Algebra
\DeclareMathOperator{\id}{id}
\DeclareMathOperator{\colim}{colim}
\DeclareMathOperator{\Aut}{Aut}
\DeclareMathOperator{\End}{End}
\DeclareMathOperator{\Hom}{Hom}
\DeclareMathOperator{\Mor}{Mor}
\DeclareMathOperator{\Ob}{Ob}
\DeclareMathOperator{\Ab}{Ab}
\DeclareMathOperator{\Coh}{Coh}
\DeclareMathOperator{\QCoh}{QCoh}
\DeclareMathOperator{\ann}{ann}
\DeclareMathOperator{\img}{img}
\DeclareMathOperator{\coker}{coker}
\DeclareMathOperator{\Ext}{Ext}
\DeclareMathOperator{\Tor}{Tor}


% Algebraic Geometry
\newcommand{\A}{\mathbb A}
\newcommand{\P}{\mathbb P}
\newcommand{\V}{\mathbb V}
\newcommand{\I}{\mathbb I}
\DeclareMathOperator{\div}{div}
\DeclareMathOperator{\Spec}{Spec}
\DeclareMathOperator{\Proj}{Proj}
\DeclareMathOperator{\Frob}{Frob}
\DeclareMathOperator{\Pic}{Pic}
\DeclareMathOperator{\Weil}{Weil}


% Unclassfified
\DeclareMathOperator{\Cinf}{C^{\infty}}
\DeclareMathOperator{\Ell}{Ell}
\DeclareMathOperator{\CL}{\mathcal{CL}}
\DeclareMathOperator{\Log}{Log}
\DeclareMathOperator{\Arg}{Arg}
\DeclareMathOperator{\Res}{Res}
\DeclareMathOperator{\val}{val}

\DeclareMathOperator{\Supp}{Supp}
\DeclareMathOperator{\cl}{cl}
\DeclareMathOperator{\disc}{disc}
\DeclareMathOperator{\M}{M}
\DeclareMathOperator{\opchar}{char}
\DeclareMathOperator{\argmax}{argmax}
\DeclareMathOperator{\argmin}{argmin}
\DeclareMathOperator{\limsup}{limsup}
\DeclareMathOperator{\Ind}{Ind}
\DeclareMathOperator{\Pr}{\mathbf{Pr}}
\DeclareMathOperator{\Exp}{\mathbb{E}}
\DeclareMathOperator{\Var}{\mathbf{Var}}

\setlength{\parindent}{0pt}
\setlength{\parskip}{5pt}
\setlength{\hfuzz}{4pt}




\begin{document}


\title{Scheme Theory}

\author{Sair Shaikh}
\maketitle



\section{Schemes: }


% \section{Affine Schemes}

% \begin{definition}
% An affine scheme $X = \Spec{A}$ is an extension of a space. It consists of data at three levels: 
% \begin{itemize}
%     \item The points of the space $X$ consists of all the prime ideals of $A$.
%     \item The topology on $X$ is defined by declaring $V(f)$ closed and extending to a topology. $V(f)$ is the vanishing set of $f \in A$, i.e. the set where $f$ evaluates to $0$. We evaluate $f \in A$ at a prime ideal $\mathfrak{p}$ by mapping $f$ into $A/\mathfrak{p}$ (or the fraction field thereof). In particular, $\mathfrak{p} \in V(f) \iff f \equiv 0 \pmod{\mathfrak{p}} \iff f \in \mathfrak{p}$. \\
%     The distinguished open sets $D(f) := X\setminus V(f)$ form a base for the topology. 
%     \item The structure sheaf $\calO_X$. Assign to $D(f)$ the localization of $A$ at the multiplicative set of all functions that do not vanish on $D(f)$, i.e. functions $g$ such that $D(f) \subseteq D(g)$. In some sense, ``rational functions'' where the denominators do not vanish on $D(f)$. The restriction maps from $D(f)$ to $D(f')$ with $D(f') \subseteq D(f)$ are given by further localization.
% \end{itemize}
    
% \end{definition}

% \subsection{The set}

% \subsection{The topology}
% \begin{definition}
%     A topological space is quasicompact if every open cover has a finite subcover. This is the AG term for compact.
% \end{definition}
% \begin{exercise}{3.5.BC and 3.6.G(a)}
%     Show that $\Spec A$ is quasicompact. 
% \end{exercise}
% \begin{proof}
%     We first show that $\Spec A = \bigcup_{i \in I} D(f_i)$ is equivalent to $(\{f_i\}_{i \in I}) = A$. \\
%     Assume $\Spec A = \bigcup_{i \in I} D(f_i)$. It suffices to show $1 \in (\{f_i\}_{i \in I})$. Suppose not. Then $(\{f_i\}_{i \in I})$ is a proper ideal of $A$ and is thus contained in some maximal ideal $\mathfrak{m}$. Thus, $\mathfrak{m} \in \bigcap_{i\in I} V(f_i)$. This implies $\mathfrak{m} \not \in \bigcup_{i \in I} D(f_i)$ which is a contradiction. Conversely, assume $(\{f_i\}_{i \in I}) = A$ and let $\mathfrak p$ be some prime ideal. Since $A \not \subseteq \mathfrak p$, there is some $f_i \not \in \mathfrak p$. Thus, $\mathfrak p \in D(f_i)$ for this $i$. \\ 
%     Now, note that $\Spec A = \bigcup_{i \in I} D(f_i)$ implies $1 = \sum_{i \in I} a_if_i$ where all but finitely many $a_i$ are zero. Thus, picking these $i$ in a set $J$, we get a finite subcover:
%     \[ \Spec A = \bigcup_{i\in J} D(f_i)\]
% \end{proof}

% \begin{exercise}{3.5.E}
%     Show that $D(f) \subseteq D(g)$ if and only if $f^n \in (g)$ if and only if $g$ is invertible in $A_f$. 
% \end{exercise}
% \begin{proof}
%     $(2) \iff (3)$ is trivial as if we have $g^{-1} = \frac{g'}{f^n} \in A_f$ with $g' \in A$ then $gg' = f^n$. \\
%     $(1) \implies (2)$. Taking the complement, $V(g) \subseteq V(f)$. Thus, $f$ vanishes at all points of $\Spec(A/(g))$. Thus, $f$ is nilpotent in $A/(g)$ so $f^n \in (g)$. More algebraically, $V(g) \subseteq V(f)$ means every prime ideal that contains $g$ also contains $f$. As $\sqrt{g} = \bigcap_{\mathfrak{p}\supseteq (g)} \mathfrak{p}$ (Ex. 3.4.F), $f \in \sqrt{g}$. \\
%     $(2) \implies (1)$. Same as above. $f^n \in (g)$ implies $f \in \sqrt{(g)}$ implies $V(g) \subseteq V(f)$ implies $D(f) \subseteq D(g)$.  
% \end{proof}


% \subsection{The structure sheaf}

% \begin{exercise}{4.1.A}
%     The natural map $A_f \to \calO_X(D(f))$ is an isomorphism.    
% \end{exercise}
% \begin{proof}
%     Exercise 3.5.E shows that $D(f) \subset D(g)$ if and only if $\frac{1}{g} \in A_f$. Thus, the map is bijective.
% \end{proof}


\end{document}